% Generated by scripts/build_readers.py; change the reviewed sources and rebuild.
\def\EditionMode{1}
\def\EditionLabel{Corrected reader with apparatus}
\def\EditionStatus{Complete checked edition}
\documentclass[10pt,oneside]{book}
\usepackage{fontspec}
\setmainfont{Latin Modern Roman}
\setsansfont{Latin Modern Sans}
\usepackage{amsmath,amssymb,mathtools}
\usepackage{unicode-math}
\setmathfont{Latin Modern Math}
\usepackage{xcolor,tikz-cd,lua-ul,marginnote,pdflscape}
\usetikzlibrary{arrows.meta,decorations.pathmorphing,calc,shapes.geometric}
\tikzcdset{ampersand replacement=\&}
\tikzcdset{every label/.append style={text=black}}
\usepackage[a4paper,left=18mm,right=45mm,top=18mm,bottom=22mm,marginparwidth=28mm,marginparsep=7mm]{geometry}
\usepackage[unicode,colorlinks=true,urlcolor=black,linkcolor=black]{hyperref}
\usepackage{bookmark}
\setlength{\parindent}{0pt}
\setlength{\parskip}{4pt plus 1pt minus 1pt}
\setlength{\emergencystretch}{1.5em}
\raggedbottom
% Include final-line descenders inside the declared text-height rectangle.
\maxdepth=0pt
\definecolor{PSHand}{RGB}{25,58,132}
\definecolor{PSDelete}{RGB}{146,37,37}
\definecolor{PSTypedOld}{RGB}{95,95,95}
\ifnum\EditionMode=2\colorlet{blue}{black}\else\colorlet{blue}{PSHand}\fi
\RenewDocumentCommand{\textcircled}{m}{\tikz[baseline=(PSCircle.base)]\node[draw,circle,inner sep=.5pt,line width=.4pt](PSCircle){#1};}
\newif\ifPSInApparatus
\newcounter{psnote}
\newcounter{psvariant}
\newcounter{pslinkmark}
\newif\ifPSInLink
\newsavebox{\PSNoteBox}
\newdimen\PSNoteLimit
{\footnotesize\global\PSNoteLimit=3\baselineskip}
\newcommand{\CurrentScanURL}{}
\newcommand{\CurrentKey}{}
\newcommand{\CurrentScan}{}
\newcommand{\pssource}[3]{\psflushmargins\clearpage\gdef\CurrentKey{#1}\gdef\CurrentScan{#2}\gdef\CurrentScanURL{#3}\setcounter{psvariant}{0}\typeout{PS-SOURCE: #1}}
\newcommand{\psdivision}[2]{\pdfbookmark[0]{#1}{#2}\par\noindent{\large\bfseries #1\par}\medskip}
\newcommand{\pslinkmarker}[2]{\stepcounter{pslinkmark}\latelua{texio.write_nl('log','PS-LINK-SHIP|#1|\CurrentKey|\arabic{pslinkmark}|#2')}}
\newcommand{\pslink}[2]{\ifPSInLink\pslinkmarker{#1}{nested}#2\else\begingroup\PSInLinktrue\href{\CurrentScanURL}{\pslinkmarker{#1}{outer}#2}\endgroup\fi}
\newcommand{\scanpage}[1]{\ifnum\EditionMode=2\leavevmode\marginpar{\scriptsize\sffamily\pslink{scanpage}{scan #1}}\else\par\noindent{\small\sffamily\pslink{scanpage}{[scan #1]}}\par\medskip\fi}
\newcommand{\mspage}[1]{\ifnum\EditionMode=2\leavevmode\marginpar{\scriptsize\sffamily p.\,#1}\else{\small\sffamily[p.\,#1]}\fi}
\newcommand{\hand}[1]{{\ifnum\EditionMode=2\color{black}\else\color{PSHand}\fi #1}}
\let\over\relax
\NewDocumentCommand{\over}{O{} m m}{{\textcolor{PSTypedOld}{#2}\ensuremath{^{\!\nearrow}}\if\relax\detokenize{#1}\relax\hand{#3}\else #3\fi}}
\NewDocumentCommand{\inserted}{O{} m}{{\if\relax\detokenize{#1}\relax\hand{\ensuremath{\langle}#2\ensuremath{\rangle}}\else\ensuremath{\langle}#2\ensuremath{\rangle}\fi}}
\newcommand{\deleted}[1]{{\color{PSDelete}\ifmmode\text{[del.: }#1\text{]}\else[\textsf{del.:} #1]\fi}}
\newcommand{\xout}[1]{{\color{PSTypedOld}\ifmmode\text{[x: }#1\text{]}\else[\textsf{x:} #1]\fi}}
\let\PSMathUnderline\underline
\RenewDocumentCommand{\underline}{+m}{\ifmmode\PSMathUnderline{#1}\else\underLine{#1}\fi}
\NewDocumentCommand{\uline}{+m}{\underline{#1}}
\NewDocumentCommand{\uuline}{+m}{\ifmmode\PSMathUnderline{\PSMathUnderline{#1}}\else\underLine{\underLine{#1}}\fi}
\newunderlinetype\PSWaveLine{\leaders\hbox{\lower.7ex\hbox{\tikz[x=1pt,y=1pt,baseline=0pt]{\draw[line width=.3pt] (0,0) .. controls (1,1) and (2,1) .. (3,0) .. controls (4,-1) and (5,-1) .. (6,0);}}}}
\NewDocumentCommand{\uwave}{+m}{{\PSWaveLine #1}}
\newcommand{\unsure}[1]{\pslink{unsure}{[#1\ifmmode\,?\else\,?\fi]}}
\newcommand{\illegible}[1]{\ifmmode\text{\pslink{illegible}{[illeg.\ #1]}}\else\pslink{illegible}{[illeg.\ #1]}\fi}
\NewDocumentCommand{\diagram}{+m}{\ifmmode\vcenter{\hbox{#1}}\else\par\medskip{\centering #1\par}\medskip\psflushmargins\fi}
\newenvironment{diagramblock}{\par\begin{center}}{\end{center}\par}
\newcommand{\PSMathNotes}{}
\makeatletter
\NewDocumentCommand{\psplacednote}{m +m}{%
      \sbox{\PSNoteBox}{\parbox{\marginparwidth}{\footnotesize\raggedright\hand{\textsf{m#1.} #2}}}%
      \ifdim\dimexpr\ht\PSNoteBox+\dp\PSNoteBox\relax<\PSNoteLimit
        \leavevmode\marginpar{\usebox{\PSNoteBox}}%
      \else
        \par\nobreak\smallskip\hrule\nobreak\smallskip
        {\footnotesize\textsf{Marginal note m#1.}\quad\hand{#2}\par}
        \nobreak\smallskip\hrule\smallskip
      \fi}
\newcommand{\psflushmargins}{\ifmmode\else\ifinner\else\ifx\PSMathNotes\@empty\else\let\PSTempNotes\PSMathNotes\global\let\PSMathNotes\@empty\PSTempNotes\fi\fi\fi}
\NewDocumentCommand{\margin}{+m}{%
  \ifPSInApparatus\ifmmode\hand{\text{[marg.: }#1\text{]}}\else\hand{\textsf{[marg.: }#1\textsf{]}}\fi
  \else\stepcounter{psnote}%
    \ifmmode
      \text{\textsuperscript{m\arabic{psnote}}}%
      \begingroup\edef\PSQueue{\endgroup\noexpand\g@addto@macro\noexpand\PSMathNotes{\noexpand\psplacednote{\arabic{psnote}}{\unexpanded{$#1$}}}}\PSQueue
    \else
      \textsuperscript{m\arabic{psnote}}\psplacednote{\arabic{psnote}}{#1}%
    \fi
  \fi}
\AtBeginDocument{\everymath\expandafter{\the\everymath\aftergroup\psflushmargins}\everydisplay\expandafter{\the\everydisplay\aftergroup\psflushmargins}}
\makeatother
\newcommand{\psnumber}[2]{\expandafter\gdef\csname psnum@#1\endcsname{#2}}
\NewDocumentCommand{\var}{+m +m m m}{\ifnum\EditionMode=0 #1\else\ifnum\EditionMode=1\ifmmode\text{\textsuperscript{\tiny\csname psnum@#4\endcsname}}\else\textsuperscript{\tiny\csname psnum@#4\endcsname}\fi\fi #2\fi}
\newcommand{\psapparatusstart}{\par\bigskip\hrule\smallskip\PSInApparatustrue\begingroup\footnotesize\raggedright\textsf{Apparatus for scan \CurrentScan}\par}
\newcommand{\psapparatusend}{\par\endgroup\PSInApparatusfalse}
\NewDocumentCommand{\psentry}{m m m +m +m m}{\par\noindent\textbf{#2. #3.}\quad\textit{Typescript:} #4\quad\textit{Corrected:} \if\relax\detokenize{#5}\relax\textnormal{[omitted]}\else #5\fi\quad #6\par}
\newcommand{\psfrontmatter}{%
\hypersetup{pageanchor=false}\begin{titlepage}\raggedright
{\Large Alexander Grothendieck\par}\vspace{12mm}
{\Huge Pursuing Stacks\par}\medskip{\LARGE 1983\par}\vspace{12mm}
{\Large\EditionLabel\par}\vspace{8mm}
{\large\EditionStatus\par}\vfill
Université de Montpellier, Grothendieck archive, cotes 134-2--134-8.\par
Including the 26 source pages of R. Brown's translation at cote 134-2, sites 115--143.\par
\vspace*{2mm}\end{titlepage}\hypersetup{pageanchor=true}
\chapter*{Editorial statement}
AI-transcribed by Claude Opus 5.5 (Anthropic, effort high); checked against the scans, and this edition built, by GPT-6 Astra (OpenAI, ChatGPT Pro mode); prompts and records by Claude (Anthropic).

The witness consulted and not reproduced is the LaTeX edition of Mateo Carmona with the collaboration of Ulrik Buchholtz (arXiv 2111.01000 v2, 2021), extended from the 2015 edition by ``the scrivener'', CC0 1.0.

The translation in the appendix is R. Brown's, prepared with assistance from L. Breen and J-L. Loday, as his note says; the rights in it belong to its translator or his heirs.

The rights in the original work belong to its author or the author's representatives or heirs. Where we hold rights, in the AI-produced transcription, translation, corrections and typesetting, they are dedicated to the public domain under CC0 1.0. This edition is published for educational, non-profit purposes, and we claim no rights we do not hold. If you hold rights in this work and object to its publication here, please open a GitHub issue and we will take it down.

\section*{How to read this form}
\ifnum\EditionMode=0
The diplomatic form preserves typed and handwritten wording, cancellations, insertions, overwritten readings, underlining, doubts, and original page numbers. Blue marks handwriting; red labels handwritten cancellations; gray labels typewriter cancellations and underlying overwritten text. Angle brackets mark insertions; a superscript arrow separates an overwritten reading from its replacement. Typed insertions and typed overwriting remain plain. The disposition of long marginal notes is identified explicitly.
\else
The corrected text applies the author's insertions, replacements, cancellations, and recorded transpositions. It retains his English and notation, including consistent forms such as ``standart'', ``developped'', ``it's'' for ``its'', and ``A la''. An isolated slip is corrected only when specifically recorded and justified. Mathematical anomalies are preserved and discussed in the guide.
\ifnum\EditionMode=1
Blue continues to identify handwriting. A numbered call identifies each maximal editorial operation, with its full diplomatic and corrected readings and reason in the apparatus following the source page. Internal layers of a wholly replaced or deleted passage are preserved within that entry rather than duplicated as separate entries. These entries are generated from the same variants that govern the running text.
\else
This reading form omits the apparatus, cancelled material, superseded readings, insertion brackets, and color distinctions. It retains the resulting text, marginal notes, underlining, and unresolved doubts. Scan and source-page references appear in the margin.
\fi\fi

Every scan mark and remaining doubt links to its source page in the archive viewer. Doubts are editorial reading marks, not emendations. The source-page breaks identify the archival sequence; a source page may occupy more than one printed page. A partial working edition is explicitly labelled and contains only reviewed source pages.
\clearpage}

\begin{document}
\pssource{ps_134-7_035}{134-7/35}{https://grothendieck.umontpellier.fr/web/viewer.html?file=../134-7.pdf\#page=36}
\psnumber{ps_134-7_035-v0001}{1}
\psnumber{ps_134-7_035-v0002}{2}
\psnumber{ps_134-7_035-v0003}{3}
\psnumber{ps_134-7_035-v0004}{4}
\psnumber{ps_134-7_035-v0005}{5}
\psnumber{ps_134-7_035-v0006}{6}
\psnumber{ps_134-7_035-v0007}{7}
\psnumber{ps_134-7_035-v0008}{8}
\psnumber{ps_134-7_035-v0009}{9}
\psnumber{ps_134-7_035-v0010}{10}
%% BEGIN SOURCE ps_134-7_035
% Pursuing Stacks (cotes 134-2 to 134-8): archive cote 134-7, site page 35 (PDF page 36), typescript page 477.
% typescript. Draft compared: ps6.tex:1668-1740.
% AI-transcribed by Claude Opus 5.5 (claude-opus-5-5, adaptive thinking, effort high), Message Batches API, batch ps3 (msgbatch_01XdFmuVVBAcgc7EFxdtCqsj), 2026-11-12; reviewed against both supplied scan halves; see CHECK_DECISIONS.tsv.
\scanpage{134-7/35}
\mspage{477}
fact,  ring extension for homotopy types (an operation of greater interest
than ring restriction surely) is \uline{not} expressed,  in general,  by just
performing the trivial ring extension operation
\[ X_{*} \hand{\mapsto} X_{*} \otimes_{k} k' \]
on models in $\hand{\pi}_{n}(k)$,  unless we assume k' to be flat over k say  -  but
even in this case it is by no means clear a priori that the operation
above transforms weak equivalences into weak equivalences.  This is very
clearly \var{\deleted{suggested}}{}{deletion}{ps_134-7_035-v0001}\var{\inserted{shown}}{\hand{shown}}{insertion}{ps_134-7_035-v0002} by the linear analogon,  the categories $\hand{\mathcal{M}}_{n}(k)$  being
replaced by \var{\inserted[typed]{the}}{the}{typed insertion}{ps_134-7_035-v0003} categories of chain complexes in  $Ab_{k}$ say,  or by $Comp^{-}(Ab_{k})$
or the like.  \var{\xout{In}}{}{typewriter deletion}{ps_134-7_035-v0004} In order to correctly describe ground ring extension on
homotopy types,  we'll have first to take a suitable``resolution'' of $X_{*}$ ,
namely replace $X_{*}$  by some  $K_{*}$  say,  endowed with a weak equivalence
\[ K_{*} \hand{\to} X_{*} \quad , \]
and $K_{*}$ satisfying some extra assumptions.  Maybe flatness of the compo-%
nent would be enough here.  For other purposes,  we may have to use
resolutions which are even smooth (componentwise),  or which satisfy a
suitable \uline{Kan condition}  (of a type outlined on page 463),  or both.
Our expectation is that,  when we restrict to the subcategory
\[ \var{\deleted{\unsure{SK_{n}(k)}}}{}{deletion}{ps_134-7_035-v0005} \ \hand{sK\mathcal{M}_{n}(k)} \]
of the model category $\hand{\mathcal{M}}_{n}(k)$ made up with smooth Kan complexes,  \uline{that}
\uline{the category} $Hot_{n}(k)$  \uline{may be described simply in  terms of such}  sK-%
\uline{complexes} \uline{``up to h\hand{$\smile$}omotopy''},  as usual.  If this is so,  the ground ring
extension functor \var{\deleted{(S)}}{}{deletion}{ps_134-7_035-v0006}  follows trivially from a corresponding functor on
the \var{\xout{\unsure{Kss}}}{}{typewriter deletion}{ps_134-7_035-v0007} sK-model categories
\[ \text{(T)} \qquad sK\hand{\mathcal{M}}_{n}(k) \hand{\longrightarrow} sK\hand{\mathcal{M}}_{n}(k') \quad , \quad \hand{X_{*} \longmapsto X_{*} \otimes_{k} k'} \hand{,} \]
hence
\[ \text{(T')} \qquad Hot_{n}(k) \hand{\longrightarrow} Hot_{n}(k') \quad . \]

From the \var{\inserted[typed]{sections}}{sections}{typed insertion}{ps_134-7_035-v0008} functor  (R')  we get homotopy invariants  $\hand{\pi}_{i}$ for an object
in  $Hot_{n}(k)$,  but the relevant k-module structure on these is not appea-%
rant on this definition.  We have a better hold,  via linear algebra over
k,  upon  homology invariants of $X_{*}$  ,  which are k-linear objects,  and
are definitely distinct (unless $k=\underline{Z}$ )  from the homology invariants of
$X_{*}(k)$ ,  which for gene ral k are definitely of little interest it seems.
The definition of homology goes via the pointed linearization functor \hand{(M)}\hand{:}
\[ \text{(U)} \qquad \hand{\Big\{} \begin{array}{l} LH.(X_{*}) \overset{def}{=} \var{\xout{\unsure{KK}}}{}{typewriter deletion}{ps_134-7_035-v0009} \ L_{\hand{pt}}\hand{(X_{*})} \var{\inserted{\text{, viewed as an object in }D.(Ab_{k}),}}{\hand{\text{, viewed as an object in }D.(Ab_{k}),}}{insertion}{ps_134-7_035-v0010} \\ \hand{H_{i}(X_{*}) \overset{def}{=} H_{i}(LH.(X_{*})) = \pi_{i}(L_{pt}(X_{*}))} \ \hand{\text{in }Ab_{k},} \end{array} \]
\margin{35}

% Visual collation: both supplied image halves inspected; see CHECK_DECISIONS_134_7_8.tsv.

%% END SOURCE ps_134-7_035
\ifnum\EditionMode=1
\psapparatusstart
\psentry{ps_134-7_035-v0001}{1}{deletion}{\deleted{suggested}}{}{The author cancels this material by hand; it leaves the corrected running text.}
\psentry{ps_134-7_035-v0002}{2}{insertion}{\inserted{shown}}{\hand{shown}}{The handwritten insertion is incorporated at its indicated position; insertion brackets leave the running text.}
\psentry{ps_134-7_035-v0003}{3}{typed insertion}{\inserted[typed]{the}}{the}{The typed insertion is incorporated at its indicated position; insertion brackets leave the running text.}
\psentry{ps_134-7_035-v0004}{4}{typewriter deletion}{\xout{In}}{}{The typewriter strikes out this material; it leaves the corrected running text.}
\psentry{ps_134-7_035-v0005}{5}{deletion}{$\deleted{\unsure{SK_{n}(k)}}$}{}{The author cancels this material by hand; it leaves the corrected running text.}
\psentry{ps_134-7_035-v0006}{6}{deletion}{\deleted{(S)}}{}{The author cancels this material by hand; it leaves the corrected running text.}
\psentry{ps_134-7_035-v0007}{7}{typewriter deletion}{\xout{\unsure{Kss}}}{}{The typewriter strikes out this material; it leaves the corrected running text.}
\psentry{ps_134-7_035-v0008}{8}{typed insertion}{\inserted[typed]{sections}}{sections}{The typed insertion is incorporated at its indicated position; insertion brackets leave the running text.}
\psentry{ps_134-7_035-v0009}{9}{typewriter deletion}{$\xout{\unsure{KK}}$}{}{The typewriter strikes out this material; it leaves the corrected running text.}
\psentry{ps_134-7_035-v0010}{10}{insertion}{$\inserted{\text{, viewed as an object in }D.(Ab_{k}),}$}{$\hand{\text{, viewed as an object in }D.(Ab_{k}),}$}{The handwritten insertion is incorporated at its indicated position; insertion brackets leave the running text.}
\psapparatusend
\fi

\end{document}
